\documentclass[UTF8,a4paper,12pt,fontset=none]{ctexart}

\usepackage[a4paper,margin=2.2cm]{geometry}
\usepackage{amsmath}
\usepackage{array}
\usepackage{booktabs}
\usepackage{caption}
\usepackage{enumitem}
\usepackage{float}
\usepackage{fontspec}
\usepackage{graphicx}
\usepackage{hyperref}
\usepackage{listings}
\usepackage{siunitx}
\usepackage{subcaption}
\usepackage{tabularx}
\usepackage{xurl}
\usepackage{xcolor}

\hypersetup{colorlinks=true,linkcolor=blue!60!black,urlcolor=blue!60!black}
\defaultCJKfontfeatures{AutoFakeSlant=.2}
\setCJKmainfont{Songti SC}
\setCJKsansfont{Heiti SC}
\setCJKmonofont{Heiti SC}
\setmonofont{Menlo}
\setlength{\parindent}{2em}
\setlength{\parskip}{0.35em}
\emergencystretch=3em
\graphicspath{{../results/figures/}}
\captionsetup{font=small}
\sisetup{round-mode=places,round-precision=3}
\renewcommand{\figurename}{图}
\renewcommand{\tablename}{表}

\definecolor{CodeBg}{HTML}{F7F7F7}
\definecolor{CodeKeyword}{HTML}{1D4E89}
\definecolor{CodeComment}{HTML}{0B6E4F}
\definecolor{CodeString}{HTML}{C73E1D}

\lstset{
  basicstyle=\ttfamily\small,
  backgroundcolor=\color{CodeBg},
  frame=single,
  framerule=0pt,
  xleftmargin=1em,
  xrightmargin=1em,
  breaklines=true,
  columns=fullflexible,
  keywordstyle=\color{CodeKeyword}\bfseries,
  commentstyle=\color{CodeComment},
  stringstyle=\color{CodeString},
  showstringspaces=false
}

\title{\bfseries 中山大学计算机学院本科生实验报告}
\author{}
\date{}

\begin{document}

\begin{titlepage}
  \centering
  {\zihao{2}\bfseries 中山大学计算机学院本科生实验报告\par}
  \vspace{0.5em}
  {\zihao{4}（2025学年春季学期）\par}
  \vspace{2em}

  \begin{tabularx}{0.9\textwidth}{>{\bfseries}p{4cm}X}
    课程名称： & 并行程序设计与算法 \\
    实验题目： & Lab9 CUDA矩阵转置 —— 基于CUDA的并行矩阵转置与性能分析 \\
    批改人： & \\
    专业（方向）： & \\
    学号： & 23336128 \\
    姓名： & 梁力航 \\
    Email： & \\
    完成日期： & 2026年5月28日 \\
  \end{tabularx}

  \vfill
  \begin{minipage}{0.88\textwidth}
    \small
    本次实验使用CUDA实现并行矩阵转置，实现了3种不同优化级别的kernel（朴素转置、
    合并读转置、共享内存分块转置），在智算习堂平台（NVIDIA A100 40GB GPU）上测试不同矩阵规模
    （512--2048）和不同的线程块大小（8×8--32×32）对性能的影响。实验发现了一个反直觉的
    结果：在中小规模矩阵下，最简单的朴素转置反而性能最优，这揭示了GPU缓存层次和编译器
    架构适配对性能的深层影响。
  \end{minipage}
\end{titlepage}

\section{实验目的}

本次实验目标如下：
\begin{enumerate}[leftmargin=2.2em]
  \item 掌握CUDA编程的基本模型：线程层次（grid/block/thread）、内存层次（global/shared/register）；
  \item 实现多种CUDA矩阵转置kernel，理解不同访存模式对性能的影响；
  \item 通过实验分析线程块大小、矩阵规模、共享内存使用等因素对GPU程序性能的作用；
  \item 学习使用CUDA Event进行精确的kernel计时和性能分析。
\end{enumerate}

\section{实验平台与参考来源}

\begin{table}[H]
  \centering
  \caption{实验平台与工具}
  \label{tab:platform}
  \begin{tabular}{ll}
    \toprule
    项目 & 参数 \\
    \midrule
    GPU & NVIDIA A100-PCIE-40GB (GA100) \\
    GPU架构 & Ampere (SM 8.0) \\
    CUDA核心 & 6,912 \\
    显存 & 40 GB HBM2e \\
    显存带宽 & 1,555 GB/s \\
    L2 Cache & 40 MB \\
    共享内存/Block & 最大 48 KB（可配置至 100 KB） \\
    Warp大小 & 32 \\
    最大线程/Block & 1024 \\
    \midrule
    CUDA Toolkit & 平台预装（nvcc，目标架构 sm\_37）\\
    C++标准 & C++14 \\
    计算平台 & 智算习堂（中山大学GPU计算平台） \\
    \midrule
    本地开发环境 & Apple M4, macOS 15 (Sequoia) \\
    Python & 3.8（平台）/ 3.13（本地 uv 管理） \\
    \bottomrule
  \end{tabular}
\end{table}

参考来源：
\begin{itemize}
  \item NVIDIA Corporation. "NVIDIA Ampere GA102 GPU Architecture Whitepaper", 2020.
  \item NVIDIA Corporation. "CUDA C++ Programming Guide", 2025.
  \item Kirk, D. B. \& Hwu, W. W. "Programming Massively Parallel Processors", 4th Edition, 2023.
  \item 课程课件："CUDA矩阵转置" 实验说明文档.
\end{itemize}

\section{问题描述与算法设计}

\subsection{问题定义}

矩阵转置是一个经典的数据重排操作：给定 $N \times N$ 矩阵 $A$，其转置矩阵 $B = A^T$，
满足 $B[j][i] = A[i][j]$（使用行优先存储，$i, j$ 为行、列索引）。

在CPU串行实现中，矩阵转置的时间复杂度为 $O(N^2)$，仅涉及数据的复制和重排，是一个典型的
内存密集型操作。在GPU上，如何高效组织线程和内存访问模式是优化矩阵转置性能的关键。

\subsection{CUDA编程模型基础}

CUDA编程模型的核心概念包括：

\begin{enumerate}[leftmargin=2.2em]
  \item \textbf{线程层次结构}：
        \begin{itemize}
          \item \textbf{Thread}：最小执行单元，执行在CUDA core上；
          \item \textbf{Warp}：32个线程为一组，以SIMT方式同时执行。同一warp内线程执行相同指令；
          \item \textbf{Block}：多个warp组成线程块，驻留在一个SM上。同一block内线程可通过共享内存
                和 \texttt{\_\_syncthreads()} 进行协作；
          \item \textbf{Grid}：多个block组成网格，对应一次kernel启动的全部线程。
        \end{itemize}

  \item \textbf{内存层次结构}：
        \begin{itemize}
          \item \textbf{Global Memory}（全局内存）：容量最大（24GB），延迟最高（$\sim$400 cycles），
                所有线程可访问。合并访存（coalesced access）至关重要：同一warp内的线程应访问连续的
                内存地址以最大化带宽利用率；
          \item \textbf{Shared Memory}（共享内存）：on-chip SRAM，低延迟（$\sim$20 cycles），
                同一block内所有线程共享。容量有限（最大48--100 KB/block），适合数据复用和交换；
          \item \textbf{Register}（寄存器）：最快，每个线程私有。溢出到local memory会严重影响性能。
        \end{itemize}

  \item \textbf{合并访存（Memory Coalescing）}：
        当同一warp的32个线程访问连续且对齐的128字节内存段时，这些访问可以被合并为一次或少数几次
        内存事务。反之，若访问是跨步（strided）或随机的，每个线程可能触发独立的内存事务，导致
        有效带宽大幅下降。
\end{enumerate}

\subsection{三种转置 Kernel 设计}

为了系统分析访存模式对性能的影响，实验实现了三种不同优化级别的转置kernel：

\begin{table}[H]
  \centering
  \caption{三种转置 Kernel 对比}
  \label{tab:kernels}
  \begin{tabular}{lccc}
    \toprule
    Kernel & 读模式 & 写模式 & 使用共享内存 \\
    \midrule
    Kernel 1: Naive（朴素） & 非合并 & 合并 & 否 \\
    Kernel 2: Coalesced Read（合并读） & 合并 & 非合并 & 否 \\
    Kernel 3: Shared Memory（共享内存分块） & 合并 & 合并 & 是 \\
    \bottomrule
  \end{tabular}
\end{table}

\subsubsection{Kernel 1 — 朴素转置}

最直观的实现：每个线程负责一个元素的转置，直接从全局内存读取 $A[col][row]$ 写入
$B[row][col]$。在行优先存储下，读取 $A[col][row]$ 等价于 $A[col \cdot N + row]$，
对于同一warp内连续线程（相邻的row），访问地址跨步为 $N$，导致读操作完全无法合并。
N较大时（$\ge 512$），每次读取可能触发32次独立的内存事务。

\begin{lstlisting}[language=C++, caption={朴素转置 Kernel（读非合并）}]
__global__ void transposeNaive(const float *A, float *B, int N) {
    int row = blockIdx.y * blockDim.y + threadIdx.y;
    int col = blockIdx.x * blockDim.x + threadIdx.x;

    if (row < N && col < N) {
        B[row * N + col] = A[col * N + row];
    }
}
\end{lstlisting}

\subsubsection{Kernel 2 — 合并读转置}

将读取方向翻转：读 $A[row][col]$ 是合并的（同warp内连续线程访问连续地址），但写入
$B[col][row]$ 是跨步的（非合并写）。由于Ampere架构对写操作的合并要求比读操作宽松
（非合并写可通过L2 cache缓冲），这种实现通常优于Kernel 1。

\begin{lstlisting}[language=C++, caption={合并读转置 Kernel}]
__global__ void transposeCoalescedRead(const float *A, float *B, int N) {
    int row = blockIdx.y * blockDim.y + threadIdx.y;
    int col = blockIdx.x * blockDim.x + threadIdx.x;

    if (row < N && col < N) {
        B[col * N + row] = A[row * N + col];
    }
}
\end{lstlisting}

\subsubsection{Kernel 3 — 共享内存分块转置}

利用共享内存实现读写均合并的最优策略：
\begin{enumerate}[leftmargin=2.2em]
  \item 每个block将 $32 \times 32$ 的子矩阵从全局内存以合并读方式载入共享内存（tile）；
  \item 使用 \texttt{\_\_syncthreads()} 同步确保整个tile加载完成；
  \item 线程读取共享内存中"转置后"的元素，以合并写方式存入全局内存目标位置。
\end{enumerate}

关键优化：共享内存 tile 声明为 \texttt{float tile[32][33]}（+1列padding），避免
32路bank冲突。因为共享内存按 $4$ 字节 $\times 32$ bank组织，同一列的不同行落在
同一bank。当block内多线程同时访问同行时，加padding使同行元素跨步偏离32的整数倍，
消除了bank冲突。

\begin{lstlisting}[language=C++, caption={共享内存分块转置 Kernel（读写均合并）}]
#define TILE_DIM 32
__global__ void transposeSharedMem(const float *A, float *B, int N) {
    __shared__ float tile[TILE_DIM][TILE_DIM + 1];  // +1 避免bank冲突

    int x = blockIdx.x * TILE_DIM + threadIdx.x;
    int y = blockIdx.y * TILE_DIM + threadIdx.y;

    // 合并读：从全局内存载入分块到共享内存
    if (x < N && y < N)
        tile[threadIdx.y][threadIdx.x] = A[y * N + x];

    __syncthreads();

    // 计算转置后的目标坐标
    int x_out = blockIdx.y * TILE_DIM + threadIdx.x;
    int y_out = blockIdx.x * TILE_DIM + threadIdx.y;

    // 合并写：从共享内存写出到全局内存
    if (x_out < N && y_out < N)
        B[y_out * N + x_out] = tile[threadIdx.x][threadIdx.y];
}
\end{lstlisting}

\section{CUDA 并行化策略分析}

\subsection{线程组织与映射}

对于 $N \times N$ 矩阵转置，使用二维线程组织天然映射到矩阵的行列结构：

\begin{itemize}
  \item \textbf{Block维度}：$B_x \times B_y$ 个线程，覆盖矩阵的一个子块；
  \item \textbf{Grid维度}：$\lceil N / B_x \rceil \times \lceil N / B_y \rceil$ 个block，
        覆盖整个矩阵；
  \item 每个线程处理矩阵中一个元素的转置，总线程数 $\ge N^2$。
\end{itemize}

\subsection{Block 大小对 Occupancy 的影响}

Occupancy（占用率）是衡量GPU资源利用率的关键指标，定义为每个SM上活跃warp数与理论上
最大warp数的比值。对于A100（Ampere, SM 8.0），每个SM最大支持：

\begin{itemize}
  \item 2048 个活跃线程（64 warps）；
  \item 最多 32 个活跃block；
  \item 最多 164 KB 共享内存/SM。
\end{itemize}

\begin{table}[H]
  \centering
  \caption{不同Block大小的资源分析（A100, SM 8.0）}
  \label{tab:occupancy}
  \begin{tabular}{lcccc}
    \toprule
    Block大小 & 线程数/Block & 最大活跃Block/SM & 活跃线程/SM & 理论占用率 \\
    \midrule
    8×8   & 64   & 32 & 2048 & 100\% \\
    16×16 & 256  & 8  & 2048 & 100\% \\
    32×32 & 1024 & 2  & 2048 & 100\% \\
    \bottomrule
  \end{tabular}
\end{table}

\begin{itemize}
  \item 8×8：每个block线程太少，虽然可容纳16个block但总线程仅1024，SM计算资源未充分利用；
  \item 16×16：最佳占用率，256线程/block × 6 block = 1536线程，恰好填满SM；
  \item 32×32：尽管占用率仅66.7\%，但每个block内warp数量多，访存延迟隐藏能力更强，
        对内存密集型操作可能更优。
\end{itemize}

\subsection{共享内存 Bank Conflict 分析}

共享内存由32个bank组成，每个bank宽度为4字节。当同一warp内多个线程同时访问相同bank
的不同地址时，发生bank conflict，访问被串行化。

在Kernel 3中，若直接使用 \texttt{tile[32][32]}：
\begin{itemize}
  \item 载入阶段：thread $(x, y)$ 访问 \texttt{tile[y][x]}，同一warp内相邻线程
        （x递增）访问同行不同列 $\to$ 不同bank，无冲突；
  \item 转置写出阶段：thread $(x, y)$ 访问 \texttt{tile[x][y]}，同一warp内
        thread $x$ 连续但访问 \texttt{tile[x][y]} 中的 $y$ 相同而 $x$ 不同，
        步长为32 $\to$ 所有线程命中同一bank！产生32路bank conflict。
\end{itemize}

解决方案：声明 \texttt{tile[32][33]}，将列维度加padding到33（一个非32倍数的值），
破坏bank编号的取模对应关系。padding后 $x \times 33$ 对32取模的结果不再是简单的
周期性重复，从而消除了bank conflict。

\section{Hello World 实验结果}

\subsection{运行方法}

在智算习堂平台上，编译运行hello\_world.cu，指定不同的grid和block维度：

\begin{lstlisting}[language=bash]
./hello_world 2 4 4   # grid=2 blocks, block=(4,4) threads
\end{lstlisting}

该配置创建 2 个线程块，每个块包含 $4 \times 4 = 16$ 个线程，共 32 个线程并行输出。

\subsection{输出顺序观察}

多次运行后发现，线程输出的顺序并非按照block编号或thread编号的线性顺序排列。
具体表现为：

\begin{enumerate}[leftmargin=2.2em]
  \item \textbf{Block间乱序}：Block 1的线程输出有时出现在Block 0的线程之前，
        说明不同block之间没有固定的调度顺序；
  \item \textbf{Block内部分有序}：同一block内，较小threadIdx.x的线程通常
        （但不绝对）先输出。这是因为同block内的warp按顺序发射，但多个warp
        之间可能交叠；
  \item \textbf{非确定性}：每次运行输出顺序不完全相同，这是GPU线程调度器动态
        调度的结果——warp调度器在可用warp中选择就绪者执行，选择策略由硬件决定。
\end{enumerate}

\textbf{结论}：GPU线程输出顺序没有规律可依赖。在实际CUDA编程中，不应该假设任何
线程执行顺序。需要同步的场景必须显式使用 \texttt{\_\_syncthreads()}（block内）
或通过多个kernel启动（grid级别）。

\section{矩阵转置实验结果与分析}

\subsection{运行时间与带宽分析}

\begin{table}[H]
  \centering
  \caption{矩阵转置运行时间汇总（各kernel最佳block大小，3次运行取平均）}
  \label{tab:runtime}
  \begin{tabular}{lcrrr}
    \toprule
    Kernel & N & 最佳Block & 平均Kernel时间(ms) & 有效带宽(GB/s) \\
    \midrule
    Naive & 512 & 16×16 & 0.00683 & 292.6 \\
    Naive & 1024 & 16×16 & 0.01536 & 546.1 \\
    Naive & 2048 & 16×16 & 0.05044 & 658.2 \\
    \midrule
    Coalesced Read & 512 & 8×8 & 0.00850 & 259.0 \\
    Coalesced Read & 1024 & 8×8 & 0.01741 & 481.9 \\
    Coalesced Read & 2048 & 8×8 & 0.05427 & 618.3 \\
    \midrule
    Shared Mem & 512 & 32×32 & 0.00742 & 262.1 \\
    Shared Mem & 1024 & 32×32 & 0.01741 & 481.9 \\
    Shared Mem & 2048 & 32×32 & 0.05667 & 595.4 \\
    \bottomrule
  \end{tabular}
\end{table}

\begin{figure}[H]
  \centering
  \includegraphics[width=0.85\textwidth]{runtime_vs_size.png}
  \caption{Kernel时间随矩阵规模的变化（最佳block大小）}
  \label{fig:runtime}
\end{figure}

\subsection{实验结果的关键发现：Naive Kernel 最优}

一个出乎预期的结果是：\textbf{最简单的朴素转置（Naive）以16×16 block大小在几乎所有测试
配置下取得了最优性能}。这与教科书上"共享内存分块是最优策略"的经典结论相矛盾，
但也正体现了GPU性能优化的复杂性。原因分析如下：

\begin{enumerate}[leftmargin=2.2em]
  \item \textbf{矩阵过小，全部在L2 Cache中}：本次实验最大矩阵N=2048仅需 $2048^2
        \times 4\text{B} = 16\text{MB}$。A100拥有40MB L2 Cache，整个矩阵（含转置结果）
        完全放入L2。当数据在L2中时，非合并访问的惩罚远小于数据必须在HBM中读取的情况，
        因为L2 cache line已经包含了相邻数据。
  \item \textbf{Compiler mismatch}：平台CUDA Toolkit仅支持到sm\_37（Kepler架构），
        编译器生成的PTX代码经过JIT编译到Ampere原生指令时，无法充分利用A100的硬件特性
        （如更大的寄存器文件、改进的warp调度器等）。
  \item \textbf{同步开销}：Shared Memory kernel中的 \texttt{\_\_syncthreads()} 是
        block级别的屏障同步。对于微秒级的极短kernel（0.006--0.05ms），同步开销
        在总执行时间中占比显著。
  \item \textbf{16×16是最佳block大小}：256线程/block在A100的SM上恰好达到
        100\%占用率（每个SM最多2048线程，可容纳8个256线程的block）。
  \item \textbf{32×32明显最差}：无论哪种kernel，32×32 block都远差于8×8和16×16。
        这可能是因为1024线程/block在大矩阵时会因为寄存器压力降低occupancy，
        且更大的block意味着每个block处理更大的子矩阵（32×32=1024元素），
        增大了缓存未命中概率。
\end{enumerate}

\subsection{带宽利用率分析}

有效带宽的计算公式为：
\begin{equation}
  \text{BW}_{\text{eff}} = \frac{2 \times N^2 \times \text{sizeof(float)}}
  {T_{\text{kernel}}}
\end{equation}

其中 $2 \times N^2 \times 4$ 字节对应读A和写B的总数据移动量。A100的理论
峰值带宽为1,555 GB/s。

实验数据显示有效带宽随矩阵规模增大而显著提升：
\begin{itemize}
  \item N=512时，最佳有效带宽仅292.6 GB/s（理论峰值的18.8\%）；
  \item N=1024时，最佳有效带宽达到546.1 GB/s（35.1\%）；
  \item N=2048时，最佳有效带宽达到658.2 GB/s（42.3\%）。
\end{itemize}

规模越大带宽利用率越高的原因是：kernel启动开销（约数微秒）在总时间中占比随
矩阵增大而减小，使得有效带宽趋近于纯数据传输的峰值。

另外，Naive kernel在32×32 block下性能急剧恶化（N=1024时从546降至356 GB/s），
这证实了非合并访问的代价在block增大时更加严重——大block中同一warp内的线程跨度更大，
非合并访问导致的cache miss更多。

\begin{figure}[H]
  \centering
  \includegraphics[width=0.85\textwidth]{bandwidth.png}
  \caption{各Kernel的有效带宽对比}
  \label{fig:bandwidth}
\end{figure}

\subsection{Block大小对性能的影响}

\begin{table}[H]
  \centering
  \caption{不同Block大小的性能对比（N=1024, Naive kernel）}
  \label{tab:blocksize}
  \begin{tabular}{lccc}
    \toprule
    Block大小 & 线程数 & 平均Kernel时间(ms) & 有效带宽(GB/s) \\
    \midrule
    8×8   & 64   & 0.01712 & 481.9 \\
    16×16 & 256  & 0.01536 & 546.1 \\
    32×32 & 1024 & 0.02355 & 356.2 \\
    \bottomrule
  \end{tabular}
\end{table}

\begin{figure}[H]
  \centering
  \includegraphics[width=0.9\textwidth]{block_size_impact.png}
  \caption{Block大小对性能的影响（N=1024）}
  \label{fig:blocksize}
\end{figure}

性能排序：$16\times16 > 8\times8 > 32\times32$，原因分析：

\begin{itemize}
  \item \textbf{16×16（256线程）}：在A100每个SM最大2048线程的限制下，256线程/block
        可容纳8个block/SM，恰好100\%占用率。同时每block内8个warp提供足够的
        并行度来隐藏访存延迟。
  \item \textbf{8×8（64线程）}：每block仅2个warp，warp调度器在遇到访存停顿时
        缺乏备选warp来切换执行，延迟隐藏能力弱。
  \item \textbf{32×32（1024线程）}：虽每block内有32个warp提供最强延迟隐藏，
        但超过SM的2048线程限制（2048/1024=2 block/SM），占用率仅50\%。
        且新增寄存器压力和指令cache竞争反而恶化了性能。
\end{itemize}

\subsection{加速比分析}

\begin{table}[H]
  \centering
  \caption{加速比汇总（相对Naive最佳配置16×16）}
  \label{tab:speedup}
  \begin{tabular}{lccc}
    \toprule
    N & Coalesced Read最佳 & Shared Mem (32×32) \\
    \midrule
    512  & 0.80×（慢25\%） & 0.92×（慢8\%） \\
    1024 & 0.88×（慢12\%） & 0.88×（慢12\%） \\
    2048 & 0.93×（慢7\%）  & 0.89×（慢11\%） \\
    \bottomrule
  \end{tabular}
\end{table}

在本实验条件下，所有"优化"kernel的加速比均小于1.0（即相对朴素实现产生了减速）。
这一结果提醒我们：优化策略的有效性高度依赖于具体场景（问题规模、硬件架构、编译器
版本），在没有实测验证的情况下不能假定理论上的优化一定会带来实际收益。

随着矩阵规模从512增大到2048，Coalesced Read的减速幅度从25\%缩小到7\%，呈现出
追赶naive的趋势。可以预测，当矩阵规模进一步增大到超出L2 cache容量时，
共享内存分块的收益将最终体现。

\begin{figure}[H]
  \centering
  \includegraphics[width=0.85\textwidth]{speedup.png}
  \caption{优化Kernel相对于Naive的加速比}
  \label{fig:speedup}
\end{figure}

\subsection{运行时间热力图}

\begin{figure}[H]
  \centering
  \includegraphics[width=0.85\textwidth]{heatmap.png}
  \caption{运行时间热力图（N × Kernel）}
  \label{fig:heatmap}
\end{figure}

\section{验证与测试}

\subsection{正确性验证方法}

\begin{enumerate}[leftmargin=2.2em]
  \item \textbf{CPU参考对比}：对每次GPU转置结果，在CPU上执行串行转置，逐元素对比
        GPU结果与CPU结果，允许 $10^{-5}$ 的浮点误差；
  \item \textbf{Checksum一致性}：计算转置结果矩阵所有元素的累加和（checksum），
        确保同一矩阵规模N下，所有kernel实现的checksum完全相同；
  \item \textbf{边界条件检查}：验证最后一行/列的不完整tile被正确处理（当N不是
        Tile大小的整数倍时）。
\end{enumerate}

\subsection{验证结果}

所有kernel实现均输出 \texttt{correct=true}，checksum在同一N下完全一致，
表明三种CUDA实现与CPU参考实现在数值上完全等价。

\section{讨论}

\subsection{为什么 Naive Kernel 最快？——L2 Cache 的关键角色}

本次实验最出乎预期的发现是：朴素转置（Naive, 16×16 block）在所有测试配置下取得了
最优性能，而理论分析预测的"最优"共享内存分块方案反而较慢。这一反直觉现象的
核心原因在于A100的巨大L2 Cache（40 MB）。

实验中的最大矩阵N=2048仅占16MB，加上转置结果共32MB，完全容纳在40MB L2 Cache内。
在传统的教学分析中，通常假设数据位于HBM（High Bandwidth Memory），非合并访问将
触发大量冗余的HBM事务。但当数据在L2中时，单次cache line加载（32B）就可以满足
同一warp内多个线程的需求，非合并访问的性能惩罚大幅降低。

这一发现说明：\textbf{课本上的优化策略有其适用范围，现代GPU的层次化缓存系统
可以显著改变不同访存模式的相对性能}。

\subsection{Compiler Mismatch 的影响}

平台CUDA Toolkit仅支持sm\_37（Kepler）作为目标架构，而实际运行硬件是Ampere A100
（sm\_80）。编译器为Kepler生成的指令通过GPU驱动JIT编译到原生Ampere指令时，
无法利用Ampere架构的改进特性（如更高效的shared memory访问、更大的寄存器文件等）。
这可能部分解释了为什么共享内存kernel的收益未能体现。

这提醒我们：在真实工程中，应确保CUDA Toolkit版本与目标GPU架构匹配，推荐至少
使用与GPU同时代的Toolkit版本。

\subsection{矩阵规模扩展性预测}

实验中Coalesced Read相对于Naive的减速幅度从25\%（N=512）缩小到7\%（N=2048），
呈现单调下降趋势。据此预测：
\begin{itemize}
  \item 当N达到4096（矩阵64MB，开始超出L2 Cache），Coalesced Read应能追平Naive；
  \item 当N达到8192（矩阵256MB，完全依靠HBM），Shared Memory的优势将充分体现，
        预计可获得1.5--3×加速比。
\end{itemize}

遗憾的是，实验要求的N上限为2048，无法在本次实验中验证这一预测。

\subsection{Shared Memory 的实际收益条件}

共享内存分块能带来显著收益的典型场景是：
\begin{itemize}
  \item 矩阵规模足够大，数据在L2中无法完全容纳；
  \item 编译器能生成针对目标架构优化的指令；
  \item TILE\_DIM需要根据GPU的shared memory大小和SM资源合理选择；
  \item 应用场景包含多次tile内的数据复用（如矩阵乘法中的累加操作），而非仅一次
        性的转置写回。
\end{itemize}

\section{结论}

本次实验成功使用CUDA实现了矩阵转置的并行化，取得了以下成果：

\begin{enumerate}[leftmargin=2.2em]
  \item 实现了三种不同优化级别的转置kernel（朴素转置、合并读转置、共享内存分块转置），
        所有实现均通过CPU参考验证（checksum完全一致）；
  \item 在NVIDIA A100 GPU上完成了完整的benchmark测试（3种矩阵规模 × 7种block/
        kernel组合 × 3次重复），收集了系统性的性能数据；
  \item 实验发现了一个反直觉的结果：在矩阵规模512--2048范围内，\textbf{最简单的朴素
        转置以16×16 block取得最优性能}，有效带宽最高达658 GB/s；
  \item 分析揭示了这一现象的根本原因——A100的40MB L2 Cache使非合并访问的
        惩罚大幅降低，同时编译器架构不匹配（sm\_37→sm\_80 JIT）限制了共享内存优化的收益；
  \item 确认了16×16 block（256线程）在A100上是最优配置，而32×32 block（1024线程）
        在所有场景下性能最差，说明了选择适当block大小的重要性；
  \item 预测了随着矩阵规模增大，共享内存分块的优势将逐步体现，为后续更大规模
        实验提供了方向。
\end{enumerate}

本次实验不仅掌握了CUDA编程的核心概念（线程组织、内存层次、合并访存、共享内存分块），
更重要的是学到了一条工程原则：\textbf{永远用实测数据验证理论优化，在GPU编程中
"显然更优"的策略可能因为缓存、编译器、硬件代际差异而产生完全相反的结果}。

\end{document}
